\begin{table}[htbp] \centering \caption{Summary Statistics China}\label{tab:summary_stats_CHN_unb_2waves} \begin{threeparttable} \begin{tabular}{l cccc} \toprule
 &  &  &  & $t$-test \\
\midrule
Location &  &   54.2\% &   45.8\% &  \\  \addlinespace \addlinespace
Log Consumption & 10.17 & 10.00 & 10.40 & -0.40*** \\
 & (0.88) & (0.87) & (0.85) &  \\  \addlinespace
Log Income & 8.71 & 8.33 & 9.17 & -0.84*** \\
 & (1.64) & (1.69) & (1.45) &  \\  \addlinespace
Female & 0.51 & 0.50 & 0.51 & -0.01** \\
 & (0.50) & (0.50) & (0.50) &  \\  \addlinespace
Age (years) & 43.34 & 43.08 & 43.69 & -0.61*** \\
 & (16.80) & (16.76) & (16.85) &  \\  \addlinespace
Education (years) & 7.54 & 6.43 & 9.00 & -2.58*** \\
 & (4.94) & (4.75) & (4.79) &  \\  \addlinespace
Household Size & 4.36 & 4.67 & 3.96 & 0.71*** \\
 & (1.85) & (1.89) & (1.71) &  \\  \addlinespace


\cmidrule{2-5}
Observations & 109,535 & 59,354 & 50,181 &  \\  \addlinespace
Individuals & 34,746 &  &  &  \\
Non-switchers &   91.2\% &  &  &  \\  \addlinespace
\bottomrule \end{tabular} \begin{tablenotes}[flushleft] \footnotesize \item{Summary statistics for China for the unbalanced panel across all waves. Source: China survey. The table reports means and standard deviations (in parentheses) based on individual-year pairs. See section 3 for further details. All variables have the same number of observations, except for income, which is missing for some observations. Income has 49,191 observations.} \end{tablenotes} \end{threeparttable} \end{table}